% ======================================================================================
% Nom     : reports/latex/sections/03_donnees.tex
% Rôle    : Section présentant le jeu de données utilisé.
% Auteur  : Maxime BRONNY
% Version : V1
% Cadre   : Réalisé dans le cadre du cours Techniques d'apprentissage artificiel
%           Master 1 Informatique Big Data
% Usage   : Fichier inclus dans main.tex lors de la compilation du rapport.
% ======================================================================================

\section{Présentation des données}
\label{sec:donnees}

\subsection{Source et vue d'ensemble}

Nous utilisons le \emph{Credit Risk Dataset} publié sur la plateforme Kaggle
\cite{kaggle_credit_risk} (fichier \texttt{data/raw/credit\_risk\_dataset.csv}).
Il contient \textbf{32\,581 observations}, chacune décrivant un emprunteur et le
prêt qui lui a été accordé, au travers de \textbf{11 variables explicatives} et
d'une variable cible binaire.

\subsection{Description des variables}

Le tableau~\ref{tab:variables} récapitule les variables du jeu de données.

\begin{table}[H]
  \centering
  \caption{Variables du \emph{Credit Risk Dataset}.}
  \label{tab:variables}
  \small
  \begin{tabular}{llp{7.2cm}}
    \toprule
    \textbf{Variable} & \textbf{Type} & \textbf{Description} \\
    \midrule
    \texttt{person\_age} & numérique & âge de l'emprunteur (années) \\
    \texttt{person\_income} & numérique & revenu annuel \\
    \texttt{person\_home\_ownership} & catégorielle (4) & logement : location, propriétaire, hypothèque, autre \\
    \texttt{person\_emp\_length} & numérique & ancienneté d'emploi (années) \\
    \texttt{loan\_intent} & catégorielle (6) & objet du prêt : personnel, études, médical, entreprise, travaux, rachat de dettes \\
    \texttt{loan\_grade} & catégorielle ordinale (A-G) & note de risque attribuée au dossier (A = meilleur) \\
    \texttt{loan\_amnt} & numérique & montant du prêt \\
    \texttt{loan\_int\_rate} & numérique & taux d'intérêt (\%) \\
    \texttt{loan\_percent\_income} & numérique & part du revenu consacrée au prêt \\
    \texttt{cb\_person\_default\_on\_file} & catégorielle (O/N) & antécédent de défaut dans l'historique de crédit \\
    \texttt{cb\_person\_cred\_hist\_length} & numérique & longueur de l'historique de crédit (années) \\
    \midrule
    \texttt{loan\_status} & \textbf{cible} (0/1) & 1 = défaut de paiement, 0 = remboursement \\
    \bottomrule
  \end{tabular}
\end{table}

\subsection{Variable cible et déséquilibre}

La cible \texttt{loan\_status} vaut $1$ pour $21{,}8\,\%$ des observations : le
jeu de données est déséquilibré, dans des proportions toutefois courantes pour
ce type de problème (figure~\ref{fig:target-distribution}). Ce déséquilibre
guide le choix des métriques (section~\ref{sec:methodologie}) et justifie le
découpage stratifié des données (section~\ref{sec:pretraitement}).

\begin{figure}[H]
  \centering
  \includegraphics[width=0.72\textwidth]{target_distribution.png}
  \caption{Distribution de la variable cible \texttt{loan\_status} :
    le jeu de données est déséquilibré ($21{,}8\,\%$ de défauts).}
  \label{fig:target-distribution}
\end{figure}

\subsection{Problèmes identifiés lors de l'exploration}

L'analyse exploratoire (figures et statistiques disponibles dans
\texttt{data/stats/}) a mis en évidence trois problèmes à traiter avant
l'apprentissage :

\begin{itemize}
  \item \textbf{valeurs manquantes} : deux variables en contiennent,
    \texttt{loan\_int\_rate} ($9{,}6\,\%$ des lignes) et
    \texttt{person\_emp\_length} ($2{,}75\,\%$) ;
  \item \textbf{valeurs aberrantes} : quelques lignes contiennent des valeurs
    physiquement impossibles, comme des âges supérieurs à 120 ans ou une
    ancienneté d'emploi de 123 ans, manifestement des erreurs de saisie ;
  \item \textbf{déséquilibre des classes} : $78/22$, discuté ci-dessus.
\end{itemize}

L'exploration montre par ailleurs des corrélations attendues : le taux de
défaut croît fortement avec le grade du prêt (de A vers G) et avec la part du
revenu consacrée au remboursement (\texttt{loan\_percent\_income}), deux
variables dont on peut anticiper qu'elles joueront un rôle important dans les
modèles. Aucune colonne n'a été jugée inutile : les 11 variables explicatives
sont conservées.
